% !TEX root = ../main.tex
\section{Diskussion}
\label{sec:diskussion}

\paragraph{Übereinstimmung mit der Theorie.} Bei Grundfrequenztaktung stimmen die gemessenen Spannungs- und Stromoberschwingungen sehr gut mit \eqref{eq:rechteck} und \eqref{eq:In} überein (Tabelle~\ref{tab:611}). Auch die Abhängigkeit $P\propto\sin\Th$ und das Minimum von $Q$ nahe $U_{\mathrm{WR1}}\approx\UN$ werden bestätigt. Der Grundschwingungsstrom liess sich nicht auf null bringen; $I_1=\SI{0.827}{\ampere}$ entspricht rund \SI{18}{\volt} Differenz zwischen $U_{\mathrm{WR1}}$ und $\UN$ und ist im Vergleich zu den Oberschwingungsströmen klein.

\paragraph{Winkel-Offset.} Die Wirkleistung wird nicht bei $\Th=0$, sondern bei $\Th\approx\SI{4.7}{\degree}$ null. Mögliche Ursachen sind die Synchronisation, Verluste in Wechselrichter, Drossel und Transformator oder Messtoleranzen. Eine Überprüfung (z.\,B.\ Messung der Phasenlage von $\uN$ und $\uWR$ am Oszilloskop) wurde nicht durchgeführt.

\paragraph{Blindleistung.} Das am PM\,3000 abgelesene $Q$ enthält die Verzerrungsblindleistung und erreicht deshalb nie null. Sein Minimum ($\Udc\approx250$ bis \SI{275}{\volt}) liegt nahe der Stelle, an der $U_{\mathrm{WR1}}=\UN$ ist; die Aufgabe verlangt jedoch $Q_1$. Da $I_1$ und $P_1$ nicht notiert wurden, konnte $Q_1$ nicht angegeben werden.

\paragraph{Lastpunkt.} Der in der Anleitung genannte Lastpunkt $P=\SI{1.5}{\kilo\watt}$, $Q_1=0$ wurde nicht erreicht; der Vergleich der Pulsmuster fand bei \SI{1.15}{\kilo\watt} statt. Weil alle Muster am selben Lastpunkt gemessen wurden, bleibt der Vergleich untereinander gültig.

\paragraph{Ausnutzung der Zwischenkreisspannung.} Für denselben Lastpunkt braucht das Trägerverfahren mit $\Udc\approx\SI{360}{\volt}$ deutlich mehr Zwischenkreisspannung als «3.\ und 5.\ eliminiert» mit \SI{282}{\volt}. Die Trägerverfahren nutzen die Gleichspannung also schlechter aus; das ist ein Nachteil, den die Theorie mit den kleineren Strom-Oberschwingungen bei hoher Schaltfrequenz aufwiegt.

\paragraph{Bestes Muster.} Nach den eigenen Messwerten hat «3.\ und 5.\ eliminiert» die kleinsten Stromoberschwingungen ($I_3=\SI{9}{\percent}$, $I_5=\SI{7}{\percent}$). Die Theorie würde bei hoher Schaltfrequenz («Trägerverfahren fein») Vorteile erwarten, weil die Oberschwingungen dann bei hohen Ordnungen liegen und von der Induktivität stark gedämpft werden. Die Messwerte zeigen das nicht: Für «fein» sind $U_3=\SI{20}{\percent}$ und $I_5=\SI{30}{\percent}$ unerwartet hoch. Mögliche Ursachen sind Übermodulation (Grundschwingung am Rand des Aussteuerbereichs), Ablese- oder Notierfehler oder ein unklarer Prozentbezug (Vermutungen). Ein eindeutiger Vergleich bräuchte $\mathrm{THD}_I$ aller Muster.
\paragraph{Unstimmigkeiten in den Messdaten.} Folgende Punkte bleiben offen und sind im Bericht als \texttt{todo}-Boxen markiert:
\begin{itemize}
  \item $Q$ bei $\Udc=\SI{350}{\volt}$ in Tabelle~\ref{tab:udc} fehlt.
  \item $P_1$, $Q_1$ und $I_3$ in Aufgabe~6.1.2 wurden nicht aufgenommen.
  \item «3.\ und 5.\ eliminiert»: $U_5=\SI{30}{\percent}$ und $I_1=\SI{1.2}{\ampere}$ sind unplausibel.
  \item Trägerverfahren: $U_3=\SI{20}{\percent}$ bis \SI{30}{\percent} ist für ein Trägerverfahren unerwartet hoch.
  \item Aufgabe~6.2.2: $P/Q$-Tabelle, $\cos\varphi$ und $\lambda$ fehlen.
  \item Aufgabe~6.3: Spurzuordnung im Oszillogramm fehlt.
  \item Aufgabe~6.4: Zuordnung der Einstellung «$-15^\circ$, \SI{288}{\volt}» zu einem Lastzustand ist unklar.
\end{itemize}
